\documentclass[11pt, a4paper]{ctexart}

\usepackage{assignment_style}

\begin{document}

% =========================================
% P87. 2
% =========================================
\begin{problem}{P87. 2 用 Jacobi 迭代、Gauss-Seidel 迭代与 SOR 迭代（$\omega = 1.46$）求解方程组，取初值 $x^{(0)} = (1, 1, 1, 1)^{\mathrm{T}}$，精确到 $10^{-3}$。}
    \[
        \begin{pmatrix}
            2 & -1 & 0 & 0 \\
            -1 & 2 & -1 & 0 \\
            0 & -1 & 2 & -1 \\
            0 & 0 & -1 & 2
        \end{pmatrix}
        x
        =
        \begin{pmatrix}
            1 \\
            0 \\
            1 \\
            0
        \end{pmatrix}.
    \]
\end{problem}

\begin{solution}
    原方程组等价于
    \[
        \begin{cases}
            x_1 = \dfrac{1 + x_2}{2}, \\
            x_2 = \dfrac{x_1 + x_3}{2}, \\
            x_3 = \dfrac{1 + x_2 + x_4}{2}, \\
            x_4 = \dfrac{x_3}{2}.
        \end{cases}
    \]
    继续求解可得精确解
    \[
        x =
        \begin{pmatrix}
            \dfrac{6}{5} \\
            \dfrac{7}{5} \\
            \dfrac{8}{5} \\
            \dfrac{4}{5}
        \end{pmatrix}
        =
        \begin{pmatrix}
            1.2 \\
            1.4 \\
            1.6 \\
            0.8
        \end{pmatrix}.
    \]

    \textbf{(1) Jacobi 迭代}

    迭代格式为
    \[
        x^{(k+1)}
        =
        \begin{pmatrix}
            \dfrac{1 + x_2^{(k)}}{2} \\
            \dfrac{x_1^{(k)} + x_3^{(k)}}{2} \\
            \dfrac{1 + x_2^{(k)} + x_4^{(k)}}{2} \\
            \dfrac{x_3^{(k)}}{2}
        \end{pmatrix}
        =
        \begin{pmatrix}
            \dfrac{1}{2} \\
            0 \\
            \dfrac{1}{2} \\
            0
        \end{pmatrix}
        +
        \begin{pmatrix}
            0 & \dfrac{1}{2} & 0 & 0 \\
            \dfrac{1}{2} & 0 & \dfrac{1}{2} & 0 \\
            0 & \dfrac{1}{2} & 0 & \dfrac{1}{2} \\
            0 & 0 & \dfrac{1}{2} & 0
        \end{pmatrix}
        x^{(k)}.
    \]

    由 $x^{(0)} = (1,1,1,1)^{\mathrm{T}}$ 出发，得到部分迭代结果：
    \[
        \renewcommand{\arraystretch}{1.15}
        \begin{array}{c|cccc}
            k & x_1^{(k)} & x_2^{(k)} & x_3^{(k)} & x_4^{(k)} \\ \hline
            0  & 1      & 1      & 1      & 1 \\
            1  & 1      & 1      & 1.5    & 0.5 \\
            2  & 1      & 1.25   & 1.25   & 0.75 \\
            3  & 1.125  & 1.125  & 1.5    & 0.625 \\
            \cdots & \cdots & \cdots & \cdots & \cdots \\
            29 & 1.1997 & 1.3989 & 1.5996 & 0.7993
        \end{array}
    \]
    当 $k = 29$ 时已达到 $10^{-3}$ 的精度要求。

    \textbf{(2) Gauss-Seidel 迭代}

    迭代格式为
    \[
        \begin{cases}
            x_1^{(k+1)} = \dfrac{1 + x_2^{(k)}}{2}, \\
            x_2^{(k+1)} = \dfrac{x_1^{(k+1)} + x_3^{(k)}}{2}, \\
            x_3^{(k+1)} = \dfrac{1 + x_2^{(k+1)} + x_4^{(k)}}{2}, \\
            x_4^{(k+1)} = \dfrac{x_3^{(k+1)}}{2}.
        \end{cases}
    \]

    部分迭代结果为
    \[
        \renewcommand{\arraystretch}{1.15}
        \begin{array}{c|cccc}
            k & x_1^{(k)} & x_2^{(k)} & x_3^{(k)} & x_4^{(k)} \\ \hline
            0  & 1      & 1      & 1      & 1 \\
            1  & 1      & 1      & 1.5    & 0.75 \\
            2  & 1      & 1.25   & 1.5    & 0.75 \\
            3  & 1.125  & 1.3125 & 1.5312 & 0.7656 \\
            \cdots & \cdots & \cdots & \cdots & \cdots \\
            12 & 1.1986 & 1.3981 & 1.5985 & 0.7992
        \end{array}
    \]
    当 $k = 12$ 时已达到 $10^{-3}$ 的精度要求。

    \textbf{(3) SOR 迭代}

    当 $\omega = 1.46$ 时，SOR 迭代格式为
    \[
        \begin{cases}
            x_1^{(k+1)}
            = \omega \dfrac{1 + x_2^{(k)}}{2} + (1 - \omega)x_1^{(k)}, \\
            x_2^{(k+1)}
            = \omega \dfrac{x_1^{(k+1)} + x_3^{(k)}}{2} + (1 - \omega)x_2^{(k)}, \\
            x_3^{(k+1)}
            = \omega \dfrac{1 + x_2^{(k+1)} + x_4^{(k)}}{2} + (1 - \omega)x_3^{(k)}, \\
            x_4^{(k+1)}
            = \omega \dfrac{x_3^{(k+1)}}{2} + (1 - \omega)x_4^{(k)}.
        \end{cases}
    \]

    部分迭代结果为
    \[
        \renewcommand{\arraystretch}{1.15}
        \begin{array}{c|cccc}
            k & x_1^{(k)} & x_2^{(k)} & x_3^{(k)} & x_4^{(k)} \\ \hline
            0  & 1      & 1      & 1      & 1 \\
            1  & 1      & 1      & 1.73   & 0.8029 \\
            2  & 1      & 1.5329 & 1.6393 & 0.8274 \\
            3  & 1.3890 & 1.5056 & 1.6790 & 0.8450 \\
            \cdots & \cdots & \cdots & \cdots & \cdots \\
            10 & 1.2003 & 1.4004 & 1.5999 & 0.8000
        \end{array}
    \]
    当 $k = 10$ 时已达到 $10^{-3}$ 的精度要求。三种方法中，SOR 收敛最快。
\end{solution}


% =========================================
% 5
% =========================================
\begin{problem}{5. 若用 Jacobi 迭代法求解方程组}
    \[
        \begin{pmatrix}
            a & -1 & -3 \\
            -1 & a & -2 \\
            3 & -2 & a
        \end{pmatrix}
        x = b,
    \]
    讨论参数 $a$ 与收敛性的关系。
\end{problem}

\begin{solution}
    Jacobi 迭代要求对角元非零，故先设 $a \ne 0$。取
    \[
        D = \operatorname{diag}(a,a,a),
    \]
    则 Jacobi 迭代矩阵为
    \[
        M = I - D^{-1}A
        =
        \begin{pmatrix}
            0 & \dfrac{1}{a} & \dfrac{3}{a} \\
            \dfrac{1}{a} & 0 & \dfrac{2}{a} \\
            -\dfrac{3}{a} & \dfrac{2}{a} & 0
        \end{pmatrix}.
    \]
    由特征方程
    \[
        \det(\lambda I - M) = 0
    \]
    可得
    \[
        \lambda^3 + \frac{4}{a^2}\lambda = 0,
    \]
    因而三个特征值为
    \[
        \lambda_1 = 0,\qquad
        \lambda_{2,3} = \pm \frac{2}{a} i.
    \]
    所以谱半径为
    \[
        \rho(M) = \max |\lambda| = \left|\frac{2}{a}\right|.
    \]
    Jacobi 迭代收敛的充要条件是 $\rho(M) < 1$，故
    \[
        \left|\frac{2}{a}\right| < 1
        \iff |a| > 2.
    \]
    综上：
    \[
        \boxed{\text{当 } a > 2 \text{ 或 } a < -2 \text{ 时，Jacobi 迭代收敛；当 } |a| \le 2 \text{ 时发散。}}
    \]
\end{solution}


% =========================================
% 6
% =========================================
\begin{problem}{6. 证明：用 Jacobi 迭代法求解二元线性方程组}
    \[
        \begin{pmatrix}
            a_{11} & a_{12} \\
            a_{21} & a_{22}
        \end{pmatrix}
        x = b
    \]
    收敛的充要条件为
    \[
        \left| \frac{a_{12}a_{21}}{a_{11}a_{22}} \right| < 1.
    \]
\end{problem}

\begin{myproof}
    设
    \[
        A =
        \begin{pmatrix}
            a_{11} & a_{12} \\
            a_{21} & a_{22}
        \end{pmatrix},
        \qquad
        D =
        \begin{pmatrix}
            a_{11} & 0 \\
            0 & a_{22}
        \end{pmatrix}.
    \]
    Jacobi 迭代矩阵为
    \[
        M = I - D^{-1}A
        =
        \begin{pmatrix}
            0 & -\dfrac{a_{12}}{a_{11}} \\
            -\dfrac{a_{21}}{a_{22}} & 0
        \end{pmatrix}.
    \]
    由特征方程
    \[
        \det(\lambda I - M)
        =
        \det
        \begin{pmatrix}
            \lambda & \dfrac{a_{12}}{a_{11}} \\
            \dfrac{a_{21}}{a_{22}} & \lambda
        \end{pmatrix}
        = \lambda^2 - \frac{a_{12}a_{21}}{a_{11}a_{22}}
        = 0
    \]
    得
    \[
        \lambda_{1,2}
        = \pm \sqrt{\frac{a_{12}a_{21}}{a_{11}a_{22}}}.
    \]
    因此
    \[
        \rho(M)
        = \max |\lambda|
        = \sqrt{\left| \frac{a_{12}a_{21}}{a_{11}a_{22}} \right|}.
    \]
    Jacobi 迭代收敛的充要条件是 $\rho(M) < 1$，故
    \[
        \sqrt{\left| \frac{a_{12}a_{21}}{a_{11}a_{22}} \right|} < 1
        \iff
        \left| \frac{a_{12}a_{21}}{a_{11}a_{22}} \right| < 1.
    \]
    结论得证。
\end{myproof}


% =========================================
% 7
% =========================================
\begin{problem}{7. 设线性方程组}
    \[
        \begin{pmatrix}
            10 & 4 & 4 \\
            4 & 10 & 8 \\
            4 & 8 & 10
        \end{pmatrix}
        x
        =
        \begin{pmatrix}
            13 \\
            11 \\
            25
        \end{pmatrix}.
    \]
    \begin{enumerate}[1.]
        \item 写出 Jacobi 迭代、Gauss-Seidel 迭代与 SOR 迭代（$\omega = 1.35$）的计算公式和迭代矩阵；
        \item 判断三种迭代方法的收敛性。
    \end{enumerate}
\end{problem}

\begin{solution}
    设
    \[
        A =
        \begin{pmatrix}
            10 & 4 & 4 \\
            4 & 10 & 8 \\
            4 & 8 & 10
        \end{pmatrix},
        \qquad
        b =
        \begin{pmatrix}
            13 \\
            11 \\
            25
        \end{pmatrix}.
    \]

    \textbf{(1) Jacobi 迭代}

    计算公式为
    \[
        \begin{cases}
            x_1^{(k+1)} = \dfrac{13 - 4x_2^{(k)} - 4x_3^{(k)}}{10}, \\
            x_2^{(k+1)} = \dfrac{11 - 4x_1^{(k)} - 8x_3^{(k)}}{10}, \\
            x_3^{(k+1)} = \dfrac{25 - 4x_1^{(k)} - 8x_2^{(k)}}{10}.
        \end{cases}
    \]
    写成矩阵形式
    \[
        x^{(k+1)} = M_1 x^{(k)} + g_1,
    \]
    其中
    \[
        M_1
        =
        \begin{pmatrix}
            0 & -\dfrac{2}{5} & -\dfrac{2}{5} \\
            -\dfrac{2}{5} & 0 & -\dfrac{4}{5} \\
            -\dfrac{2}{5} & -\dfrac{4}{5} & 0
        \end{pmatrix},
        \qquad
        g_1
        =
        \begin{pmatrix}
            \dfrac{13}{10} \\
            \dfrac{11}{10} \\
            \dfrac{5}{2}
        \end{pmatrix}.
    \]

    \textbf{(2) Gauss-Seidel 迭代}

    计算公式为
    \[
        \begin{cases}
            x_1^{(k+1)} = \dfrac{13 - 4x_2^{(k)} - 4x_3^{(k)}}{10}, \\
            x_2^{(k+1)} = \dfrac{11 - 4x_1^{(k+1)} - 8x_3^{(k)}}{10}, \\
            x_3^{(k+1)} = \dfrac{25 - 4x_1^{(k+1)} - 8x_2^{(k+1)}}{10}.
        \end{cases}
    \]
    记 $A = D + L + U$，则
    \[
        x^{(k+1)} = M_2 x^{(k)} + g_2,
        \qquad
        M_2 = -(D+L)^{-1}U,
        \qquad
        g_2 = (D+L)^{-1}b.
    \]
    计算得
    \[
        M_2
        =
        \begin{pmatrix}
            0 & -\dfrac{2}{5} & -\dfrac{2}{5} \\
            0 & \dfrac{4}{25} & -\dfrac{16}{25} \\
            0 & \dfrac{4}{125} & \dfrac{84}{125}
        \end{pmatrix},
        \qquad
        g_2
        =
        \begin{pmatrix}
            \dfrac{13}{10} \\
            \dfrac{29}{50} \\
            \dfrac{379}{250}
        \end{pmatrix}.
    \]

    \textbf{(3) SOR 迭代（$\omega = 1.35$）}

    计算公式为
    \[
        \begin{cases}
            x_1^{(k+1)}
            =
            \dfrac{\omega(13 - 4x_2^{(k)} - 4x_3^{(k)})}{10}
            + (1-\omega)x_1^{(k)}, \\
            x_2^{(k+1)}
            =
            \dfrac{\omega(11 - 4x_1^{(k+1)} - 8x_3^{(k)})}{10}
            + (1-\omega)x_2^{(k)}, \\
            x_3^{(k+1)}
            =
            \dfrac{\omega(25 - 4x_1^{(k+1)} - 8x_2^{(k+1)})}{10}
            + (1-\omega)x_3^{(k)}.
        \end{cases}
    \]
    其迭代矩阵可写成
    \[
        x^{(k+1)} = M_3 x^{(k)} + g_3,
        \qquad
        M_3 = (D+\omega L)^{-1}\bigl[(1-\omega)D-\omega U\bigr],
        \qquad
        g_3 = \omega (D+\omega L)^{-1}b.
    \]
    代入 $\omega = 1.35$，可得
    \[
        M_3
        \approx
        \begin{pmatrix}
            -0.35 & -0.54 & -0.54 \\
            0.189 & -0.0584 & -0.7884 \\
            -0.01512 & 0.354672 & 0.793072
        \end{pmatrix},
        \qquad
        g_3
        \approx
        \begin{pmatrix}
            1.755 \\
            0.5373 \\
            1.847016
        \end{pmatrix}.
    \]

    \textbf{(4) 收敛性判断}

    计算三个迭代矩阵的谱半径，得到
    \[
        \rho(M_1) \approx 1.0928,\qquad
        \rho(M_2) \approx 0.6283,\qquad
        \rho(M_3) \approx 0.4463.
    \]
    因而
    \[
        \boxed{\text{Jacobi 迭代发散，而 Gauss-Seidel 迭代和 SOR 迭代都收敛。}}
    \]
    且在这里 SOR 迭代的谱半径最小，因此收敛速度最快。
\end{solution}


% =========================================
% 9
% =========================================
\begin{problem}{9. 设线性方程组}
    \[
        \begin{pmatrix}
            1 & a & a \\
            4a & 1 & 0 \\
            a & 0 & 1
        \end{pmatrix}
        x
        =
        \begin{pmatrix}
            b_1 \\
            b_2 \\
            b_3
        \end{pmatrix}.
    \]
    \begin{enumerate}[1.]
        \item 写出 Jacobi 迭代和 Gauss-Seidel 迭代的计算公式；
        \item 用迭代收敛的充要条件给出使这两种迭代都收敛的 $a$ 的取值范围。
    \end{enumerate}
\end{problem}

\begin{solution}
    \textbf{(1) 迭代公式}

    Jacobi 迭代为
    \[
        \begin{cases}
            x_1^{(k+1)} = b_1 - a x_2^{(k)} - a x_3^{(k)}, \\
            x_2^{(k+1)} = b_2 - 4a x_1^{(k)}, \\
            x_3^{(k+1)} = b_3 - a x_1^{(k)}.
        \end{cases}
    \]

    Gauss-Seidel 迭代为
    \[
        \begin{cases}
            x_1^{(k+1)} = b_1 - a x_2^{(k)} - a x_3^{(k)}, \\
            x_2^{(k+1)} = b_2 - 4a x_1^{(k+1)}, \\
            x_3^{(k+1)} = b_3 - a x_1^{(k+1)}.
        \end{cases}
    \]

    \textbf{(2) 收敛区间}

    对 Jacobi 迭代，有 $D = I$，故
    \[
        M_J = I - D^{-1}A = I - A
        =
        \begin{pmatrix}
            0 & -a & -a \\
            -4a & 0 & 0 \\
            -a & 0 & 0
        \end{pmatrix}.
    \]
    特征方程为
    \[
        \det(\lambda I - M_J)
        =
        \det
        \begin{pmatrix}
            \lambda & a & a \\
            4a & \lambda & 0 \\
            a & 0 & \lambda
        \end{pmatrix}
        =
        \lambda^3 - 5a^2 \lambda
        = 0.
    \]
    因而
    \[
        \lambda_1 = 0,\qquad
        \lambda_{2,3} = \pm \sqrt{5}\,a,
    \]
    所以
    \[
        \rho(M_J) = \sqrt{5}\,|a|.
    \]
    Jacobi 收敛当且仅当
    \[
        \sqrt{5}\,|a| < 1
        \iff
        |a| < \frac{1}{\sqrt{5}}
        = \frac{\sqrt{5}}{5}.
    \]

    对 Gauss-Seidel 迭代，取
    \[
        L =
        \begin{pmatrix}
            0 & 0 & 0 \\
            4a & 0 & 0 \\
            a & 0 & 0
        \end{pmatrix},
        \qquad
        U =
        \begin{pmatrix}
            0 & a & a \\
            0 & 0 & 0 \\
            0 & 0 & 0
        \end{pmatrix},
    \]
    则
    \[
        M_G = -(D+L)^{-1}U
        =
        \begin{pmatrix}
            0 & -a & -a \\
            0 & 4a^2 & 4a^2 \\
            0 & a^2 & a^2
        \end{pmatrix}.
    \]
    其特征值为
    \[
        0,\ 0,\ 5a^2,
    \]
    因而
    \[
        \rho(M_G) = 5a^2.
    \]
    Gauss-Seidel 收敛当且仅当
    \[
        5a^2 < 1
        \iff
        |a| < \frac{\sqrt{5}}{5}.
    \]

    故使两种方法都收敛的参数范围为
    \[
        \boxed{-\frac{\sqrt{5}}{5} < a < \frac{\sqrt{5}}{5}.}
    \]
\end{solution}


% =========================================
% 12
% =========================================
\begin{problem}{12. 证明：对称矩阵}
    \[
        A =
        \begin{pmatrix}
            1 & a & a \\
            a & 1 & a \\
            a & a & 1
        \end{pmatrix}
    \]
    当 $-\dfrac{1}{2} < a < 1$ 时正定，但只有当 $-\dfrac{1}{2} < a < \dfrac{1}{2}$ 时，Jacobi 迭代求解 $Ax=b$ 才收敛。
\end{problem}

\begin{myproof}
    \textbf{(1) 判定正定性}

    由 Sylvester 判据，只需考察顺序主子式：
    \[
        \Delta_1 = 1 > 0,
    \]
    \[
        \Delta_2 =
        \begin{vmatrix}
            1 & a \\
            a & 1
        \end{vmatrix}
        = 1 - a^2 > 0
        \iff -1 < a < 1,
    \]
    \[
        \Delta_3 =
        \begin{vmatrix}
            1 & a & a \\
            a & 1 & a \\
            a & a & 1
        \end{vmatrix}
        = (1-a)^2(1+2a) > 0
        \iff -\frac{1}{2} < a < 1.
    \]
    综合可得
    \[
        \boxed{A \text{ 在 } -\frac{1}{2} < a < 1 \text{ 时正定。}}
    \]

    \textbf{(2) Jacobi 迭代的收敛性}

    由于 $D = I$，Jacobi 迭代矩阵为
    \[
        M = I - D^{-1}A = I - A
        =
        \begin{pmatrix}
            0 & -a & -a \\
            -a & 0 & -a \\
            -a & -a & 0
        \end{pmatrix}.
    \]
    该矩阵的特征值为
    \[
        -2a,\quad a,\quad a.
    \]
    因而谱半径为
    \[
        \rho(M) = \max\{2|a|, |a|\} = 2|a|.
    \]
    Jacobi 迭代收敛的充要条件是
    \[
        \rho(M) < 1
        \iff
        2|a| < 1
        \iff
        |a| < \frac{1}{2}.
    \]
    即
    \[
        \boxed{-\frac{1}{2} < a < \frac{1}{2}.}
    \]

    因此，虽然矩阵在区间 $-\dfrac{1}{2} < a < 1$ 上正定，但 Jacobi 迭代只有在更小的区间
    \[
        -\frac{1}{2} < a < \frac{1}{2}
    \]
    内才收敛。
\end{myproof}

\end{document}
